\section{Analysis}
\label{sec:analysis}

Having established \emph{where} the animals are lost, we ask \emph{which} animals and
\emph{why}.

\subsection{Where the 36 missed animals go}

\begin{table}[t]
\centering\small
\begin{tabular}{lrrrr}
\toprule
Outcome & $n$ & Median box & GT frames & Cand.\ conf \\
\midrule
Counted & 47 & 38.1\,px & 102 & 0.76 \\
\textbf{Rejected by the rule} & \textbf{27} & 28.0\,px & 47 & 0.52 \\
Never detected & 9 & 19.9\,px & 47 & --- \\
\bottomrule
\end{tabular}
\caption{\textbf{Fate of the 83 unseen animals.} The detector finds 74 of 83; the
confirmation rule then discards 27 of them. Missed animals are smaller and shorter-lived
than counted ones.}
\label{tab:missed}
\end{table}

\Cref{tab:missed} decomposes the outcome per animal. The detector never fires on 9 animals
(median 20\,px --- at the sensor's limit). The rule then rejects 27 animals the detector
\emph{did} find. Of those 27, \textbf{26 failed the confidence test} and only one failed
purely on track length: they are faint, distant animals whose detector confidence never
reaches 0.65, and a scalar threshold cannot separate them from noise. Those 26 animals are
the entire remaining opportunity, and they are what a learned confirmer would have to
recover.

We verified by direct inspection that none of the 36 are annotation errors: each was
rendered with its ground-truth box at three points in its track and checked. All are real
animals.

\subsection{The detector fails at both scale extremes}

\begin{figure}[t]
  \centering
  \includegraphics[width=\linewidth]{figures/fig4_scale.pdf}
  \caption{\textbf{A hole in the training distribution.} The training split contains no box
  above 96\,px. Three missed animals sit at 104, 149 and 157\,px --- one of them never
  detected at all. A detector that finds 27\,px animals does not miss a 149\,px one unless it
  has never seen one.}
  \label{fig:scale}
\end{figure}

That inspection produced an unexpected finding. Three of the missed animals are not small
but very \emph{large}: 104, 149 and 157\,px, against a corpus median of 27. One of them, at
149\,px, was never detected at all --- implausible for a detector that routinely finds
27-pixel animals.

\Cref{fig:scale} gives the likely explanation. The largest box in the training split is
95.6\,px and the split contains \emph{zero} boxes at or above 100\,px, so all three
animals sit outside the annotated size distribution. We state this carefully: scale
augmentation does present enlarged copies during training, so the detector is not strictly
naive to large objects --- but enlargement is sampled rarely and mosaic augmentation shrinks
objects at the same time, so large presentations are uncommon. The mirror image of the
small-object problem the field concentrates on is nonetheless visible here: animals close
to the road are rare in our data, and the detector performs worst exactly where examples
are scarcest.

The remedy is less straightforward than it first appears, and the attempt is instructive.
Widening scale augmentation to cover 100--200\,px collapsed training to zero mAP within
five epochs, because the standard implementation is symmetric: enlarging the sampling range
also shrinks a 27\,px animal to sub-pixel size, and on a corpus that is 71\% small objects
those examples are unlearnable. Enlarging the upper tail without destroying the lower one
requires asymmetric augmentation that common frameworks do not expose. The loss is
concentrated --- the single video containing these animals counts 3 of 8, and three of its
five misses are the oversized individuals --- so this is worth roughly 4 points of our
headline figure. We report the diagnosis and the failed remedy, and note that most
small-object work does not check the upper tail at all.

\subsection{Individual re-identification at 27 pixels}
\label{sec:reid}

A natural proposal for the association stage is appearance-based individual
re-identification, keyed on body shape, thermal silhouette, relative size, and movement.
Because the annotation gives one track per individual, the corpus already contains identity
labels, and the question can be settled empirically.

We evaluate closed-set re-identification on \textbf{69 animals never seen, in 7 videos never
seen}: each animal's track is split temporally, the first half forms the gallery and the
second half the query, and we report rank-1 accuracy. Within-video is the \emph{ceiling} --- both
halves share range, pose, and thermal calibration --- so a low score here is decisive.

\begin{table}[t]
\centering\small
\begin{tabular}{lc}
\toprule
Cue & Rank-1 \\
\midrule
Chance & 0.101 \\
Mean intensity alone & 0.101 \\
Thermal silhouette alone (scale removed) & 0.435 \\
Raw thermal crop & 0.493 \\
\textbf{Box size alone ($w,h$)} & \textbf{0.638} \\
ImageNet ResNet-50 features & 0.725 \\
Geometry (size, aspect, intensity) & 0.797 \\
Trained thermal encoder (ours) & \textbf{0.899} \\
\bottomrule
\end{tabular}
\caption{\textbf{Re-identification is dominated by size, not appearance.} All cues scored on
the same 69 held-out animals. A two-number size descriptor reaches 0.638 and an ImageNet
backbone scores \emph{below} plain geometry. The trained encoder's 0.899 is largely range
matching, not identity.}
\label{tab:reid}
\end{table}

\Cref{tab:reid} is the result. A trained thermal encoder reaches rank-1 0.899, which sounds
like working re-identification --- until the decomposition. \textbf{Box size alone reaches
0.638}, and within a single video an animal's range changes slowly, so most of what looks
like identification is matching \emph{distance}. Appearance in isolation is weak: mean
intensity carries exactly chance information, silhouette with scale removed reaches 0.435,
and a pretrained CNN lands below plain geometry.

We therefore describe our system as performing \emph{track association} --- preventing double
counting within a survey --- and make no individual-identification claim. At a median
$29\times24$ pixels there is not enough detail on the animal to key an identity to. We
report this as a measured limit rather than a discarded experiment, because the negative
result is the useful part: it tells practitioners where the sensor resolution boundary lies.

Adding appearance to the \emph{association} stage does help
(\cref{tab:pools}, pool G) --- it merges fragments of one animal within a video, which is a
much easier problem than identifying an individual across sightings.

\subsection{Consolidated ablations}

\begin{table}[t]
\centering\small
\setlength{\tabcolsep}{3pt}
\begin{tabular}{llc}
\toprule
Ablation & Range tested & Effect \\
\midrule
CLAHE normalization & on / off & 0.519 vs 0.299 mAP50 \\
Input resolution & 640 / 1280 & architecture-dependent \\
Detector architecture & 12 models & 0.365--0.510 AP$_{50}$ \\
Matching criterion & 4 rules & re-ranks 8/12 models \\
Detector confidence & 0.10/0.05/0.02 & 0 extra animals \\
Track-init threshold & 0.15 / 0.05 & $+5$ \primaryq \\
Orphan recovery & on / off & $+17$ animals \\
NMS IoU & 0.50 / 0.90 & $+4$ \primaryq, $-14$ \countedq \\
Appearance ReID & on / off & $+7$ \primaryq, $-3$ \countedq \\
Confirmer capacity & 3--60k params & monotone degradation \\
ReID cue & 7 variants & size dominates \\
Track splitting & threshold sweep & $+3$ for 4k cands \\
\bottomrule
\end{tabular}
\caption{\textbf{Twelve ablation families.} Every stage of the pipeline was varied
independently. The consistent pattern: interventions that improve the candidate pool do not
improve the count.}
\label{tab:ablations}
\end{table}

\Cref{tab:ablations} consolidates every ablation. One entry is worth isolating: lowering detector confidence from 0.10 to 0.05 added
3,807 candidates and recovered \emph{zero} additional animals. That single measurement
ruled out the detector threshold and pointed us at the tracker's initialization gate, which
turned out to be the real constraint.

\subsection{Qualitative results}

\begin{figure}[t]
  \centering
  \begin{subfigure}{\linewidth}
    \includegraphics[width=\linewidth]{figures/qualitative/NShelbyRdblue_SHB_12.11.2025_LS__f8544_6deer.jpg}
  \end{subfigure}\\[2pt]
  \begin{subfigure}{\linewidth}
    \includegraphics[width=\linewidth]{figures/qualitative/GolfDr_SHB_12.11.2025__f3652_5deer.jpg}
  \end{subfigure}
  \caption{\textbf{Group scenes on held-out video.} Animals standing close together are
  resolved as distinct tracked individuals rather than merged. Labels give track ID and
  detector confidence; colour is stable per track. Both frames are from videos the detector
  never saw.}
  \label{fig:qualitative}
\end{figure}

\Cref{fig:qualitative} shows representative held-out frames. Group scenes --- the hardest
case for counting, and the one where per-frame detection metrics are least informative ---
are resolved into distinct individuals.
